\documentclass[11pt]{article}
\usepackage[UTF8]{ctex}
\usepackage[a4paper]{geometry}
\geometry{left=2.0cm,right=2.0cm,top=2.5cm,bottom=2.5cm}

\usepackage{comment}
\usepackage{booktabs}
\usepackage{graphicx}
\usepackage{diagbox}
\usepackage{amsmath,amsfonts,graphicx,amssymb,bm,amsthm}
%\usepackage{algorithm,algorithmicx}
\usepackage[ruled]{algorithm2e}
\usepackage[noend]{algpseudocode}
\usepackage{fancyhdr}
\usepackage{tikz}
\usepackage{graphicx}
\usetikzlibrary{arrows,automata}
\usepackage{hyperref}
\hypersetup{
	colorlinks=true,
	linkcolor=blue,
	filecolor=blue,      
	urlcolor=blue,
	citecolor=cyan,
}			

\setlength{\headheight}{14pt}
\setlength{\parindent}{0 in}
\setlength{\parskip}{0.5 em}

\newtheorem{theorem}{Theorem}
\newtheorem{lemma}[theorem]{Lemma}
\newtheorem{proposition}[theorem]{Proposition}
\newtheorem{claim}[theorem]{Claim}
\newtheorem{corollary}[theorem]{Corollary}
\newtheorem{definition}[theorem]{Definition}
\newtheorem*{definition*}{Definition}

\newenvironment{problem}[2][Problem]{\begin{trivlist}
\item[\hskip \labelsep {\bfseries #1}\hskip \labelsep {\bfseries #2.}]}{\hfill$\blacktriangleleft$\end{trivlist}}
\newenvironment{answer}[1][Answer]{\begin{trivlist}
\item[\hskip \labelsep {\bfseries #1.}\hskip \labelsep]}{\hfill$\lhd$\end{trivlist}}

\newcommand\E{\mathbb{E}}
\newcommand\per{\mathrm{per}}


\title{Homework \#1}
\usetikzlibrary{positioning}

\begin{document}

\pagestyle{fancy}
\lhead{Peking University}
\chead{}
\rhead{Mathematical Foundations for the Information Age, 2025 Fall}

\begin{center}
    {\LARGE \bf Homework \#1}\\
    {Due: 2025-10-12 23:59 \quad$|$\quad 8 Problems, 100 Pts}\\
    {Name: Rainybunny, ID: 20120712}            % Write down your name and ID here.
\end{center}



\begin{problem}{1 (8')}
Define
\begin{align*}
    \Gamma(s) = \int_{0}^{+\infty} x^{s - 1} \mathrm{e}^{-x} \mathrm{d}x.
\end{align*}
where $s>0$. It can be proved that $\Gamma(s)$ is well-defined (You don't need to prove this).
\begin{itemize}
    \item [(1)] (4') Prove that, $\Gamma(s+1)=s\Gamma(s)$.
    \item [(2)] (4') Prove that,
    \begin{align*}
        \Gamma(s)=2\int_{0}^{+\infty}x^{2s-1}\mathrm{e}^{-x^2}\mathrm{d}x.
    \end{align*}
\end{itemize}
\end{problem}

\begin{answer} ~
\begin{itemize}
    \item [(1)]
        $$
            \begin{aligned}
                \Gamma(s) &= \int_{0}^{+\infty} x^{s-1}\mathrm{e}^{-x}\mathrm{d}x\\
                &= \frac{1}{s}\int_{0}^{+\infty} \mathrm{e}^{-x}\mathrm{d}x^s\\
                &= \frac{1}{s}\left(\left.\mathrm{e}^{-x}x^s\right|_0^{+\infty} + \Gamma(s + 1)\right)\\
                &= \frac{1}{s}\Gamma(s + 1).
            \end{aligned}
        $$
    \item [(2)]
        $$
            \begin{aligned}
                \Gamma(s) &= \int_{0}^{+\infty} (t^2)^{s-1} \mathrm{e}^{-t^2}\mathrm{d} t^2\\
                &= 2\int_{0}^{+\infty} x^{2s-1} \mathrm{e}^{-x^2} \mathrm{d}x.
            \end{aligned}
        $$
\end{itemize}
\end{answer}



\begin{problem}{2 (10')}
Define a random variable $Q \sim \chi^2(k)$ $(k\in\mathbb{Z}_+)$ if $Q=Z_1^2+\cdots+Z_k^2$ where $Z_1,\cdots,Z_k \sim \mathcal{N}(0,1)$ are independent random variables. Given two independent random variables $X\sim\chi^2(m),\;Y\sim\chi^2(n)$ $(m,n\in\mathbb{Z}_+,\;m > n)$.
\begin{itemize}
    \item [(1)] (4') Calculate the value of $\mathbb{E}(X)$.
    \item [(2)] (2') Prove that, $X + Y \sim \chi^2(m + n)$.
    \item [(3)] (4') Does $X - Y \sim \chi^2(m - n)$? Prove your result.
\end{itemize}
\end{problem}

\begin{answer} ~
\begin{itemize}
    \item [(1)]
        $$
            \mathbb{E}(X) = \sum_{i=1}^{m} \mathbb{E}(Z_i^2) = \sum_{i=1}^{m} 1 = m.
        $$
    \item [(2)]
        $$
            \begin{aligned}
                X+Y &= Z_{X1}^2+\cdots+Z_{Xm}^2+Z_{Y1}^2+\cdots+Z_{Yn}^2\\
                &= Z_1^2+\cdots+Z_{m+n}^2\\
                &\sim \chi^2(m+n).
            \end{aligned}
        $$
    \item [(3)] No. $\mathbb P[X-Y<0]>0$, while $\mathbb P[Z<0]=0$ for any $Z\sim\chi^2(k)$.
\end{itemize}
\end{answer}



\begin{problem}{3 (8')} In this problem, we will prove two basic probability inequalities. You will use these inequalities frequently in this course, so make sure you are familiar with them, including the statements and conditions.
\begin{itemize}
    \item [(1)] (4') (Markov Inequality) Suppose $X$ is a non-negative random variable and $\mathbb{E}(X)<+\infty$. Prove that, for any $a>0$,
    \begin{align*}
        \mathbb{P}(X\geq a)\leq\frac{\mathbb{E}(X)}{a}.
    \end{align*}
    \item [(2)] (4') (Chebyshev Inequality) Suppose $X$ is a random variable with $\mathbb{E}(X)<+\infty, Var(X)<+\infty$. Prove that, for any $a>0$,
    \begin{align*}
        \mathbb{P}(|X-\mathbb{E}(X)|\geq a)\leq\frac{Var(X)}{a^2}.
    \end{align*}
\end{itemize}
\end{problem}

\begin{answer} ~
\begin{itemize}
    \item [(1)] We may assume $a>\mathbb{E}(X)$ and $X\sim\rho$. In this case,
        $$
            \begin{aligned}
                \mathbb{E}(X) &\ge \int_{a}^{+\infty} x\mathrm{d}\rho(x) \\
                &\ge a \int_a^{+\infty} \mathrm{d}\rho(x)\\
                &= a\mathbb{P}(X\ge a),
            \end{aligned}
        $$
        which gives the goal.
    \item [(2)] Similarly, let $X\sim\rho$, then
        $$
            \begin{aligned}
                \mathrm{Var}(X) &= \int_{-\infty}^{+\infty}(x-\mathbb{E}(X))^2\mathrm{d}\rho\\
                &\ge \int_{|x-\mathbb{E}(X)|\ge a}(x-\mathbb{E}(X))^2\mathrm{d}\rho\\
                &\ge a^2\int_{|x-\mathbb{E}(X)|\ge a}\mathrm{d}\rho\\
                &= a^2\mathbb{P}(|X-\mathbb{E}(X)|\ge a),
            \end{aligned}
        $$
        which gives the goal.
\end{itemize}
\end{answer}



\begin{problem}{4 (14')}
In this problem, we will prove the following Chernoff bound (Theorem \ref*{thm:chernoff}).
\begin{theorem}[Chernoff Bound]
\label{thm:chernoff}
    Suppose $X_1,\cdots,X_n$ are independently Bernoulli random variables with expectation $p\in(0,1)$. Then for any $\varepsilon\in(0,1-p)$,
    \begin{align*}
        \mathbb{P}\left(\frac{1}{n}\sum_{i=1}^n X_i\geq p+\varepsilon\right)\leq \exp\left[-nD_B(p+\varepsilon||p)\right] \leq \exp(-2n\varepsilon^2),
    \end{align*} 
where 
\begin{align*}
    D_B(p||q):=p\ln\frac{p}{q}+(1-p)\ln\frac{1-p}{1-q}.
\end{align*}
\end{theorem}
\begin{itemize}
    \item [(1)] (5') Under the conditions of Theorem \ref*{thm:chernoff}, prove that for any $t>0$,
    \begin{align*}
        \mathbb{P}\left(\sum_{i=1}^n X_i\geq n(p+\varepsilon)\right)\leq\mathbb{E}\left(\mathrm{e}^{t\sum_{i=1}^n X_i}\right)\cdot\mathrm{e}^{-nt(p+\varepsilon)}.
    \end{align*}
    \item [(2)] (9') Finish the proof of Theorem \ref*{thm:chernoff}.
\end{itemize}
\end{problem}

\begin{answer} ~
\begin{itemize}
    \item [(1)] For any $t>0$,
        $$
            \begin{aligned}
                \mathbb{P}\left(\sum_{i=1}^n X_i\ge n(p+\varepsilon)\right)
                &= \sum_{s\ge n(p+\varepsilon)} \binom{n}{s} p^s\\
                &\le \sum_{s\ge n(p+\varepsilon)} \binom{n}{s} p^s \mathrm{e}^{t(s-n(p+\varepsilon))}\\
                &\le \sum_{s} \binom{n}{s} p^s \mathrm{e}^{ts} \cdot \mathrm{e}^{-tn(p+\varepsilon)}\\
                &= \mathbb{E}\left(\mathrm{e}^{t\sum_{i=1}^{n}X_i}\right)\cdot \mathrm{e}^{-tn(p+\varepsilon)}.
            \end{aligned}
        $$
    \item [(2)] \textit{For the second inequality,}
        $$
            \begin{aligned}
                &\exp\left[-nD_B(p+\varepsilon||p)\right] \leq \exp(-2n\varepsilon^2)\\
                \Leftrightarrow & \exp\left[-n\left((p+\varepsilon)\ln\frac{p+\varepsilon}{p}+(1-p-\varepsilon)\ln\frac{1-p-\varepsilon}{1-p}\right)\right] \leq \exp(-2n\varepsilon^2)\\
                \Leftrightarrow & (p+\varepsilon)\ln\frac{p+\varepsilon}{p}+(1-p-\varepsilon)\ln\frac{1-p-\varepsilon}{1-p} \geq 2\varepsilon^2.\\
            \end{aligned}
        $$

        Let
        $$
            f(\varepsilon):=(p+\varepsilon)\ln\frac{p+\varepsilon}{p}+(1-p-\varepsilon)\ln\frac{1-p-\varepsilon}{1-p}-2\varepsilon^2\quad (\varepsilon\in(0,1-p)).
        $$
        We need to prove $f(\varepsilon)\ge 0$. Note that
        $$
            f'(\varepsilon)=\ln\left(1+\frac{\varepsilon}{p}\right)+\ln\left(1-\frac{\varepsilon}{1-p}\right)-4\varepsilon,
        $$
        and
        $$
            \begin{aligned}
                f''(\varepsilon) &= \frac{1}{p+\varepsilon}+\frac{1}{1-p-\varepsilon}-4\\
                &= \frac{1}{(p+\varepsilon)(1-p-\varepsilon)}-4\\
                &\ge 0.
            \end{aligned}
        $$

        Therefore, $f'(\varepsilon)\ge f'(0+)=0$ and $f(\varepsilon)\ge f(0+)=0$.

        \,

        \textit{For the first inequality,} by the conclusion of (1), it remains to prove that,
        there exists a $t>0$ s.t.
        $$
            g(t)^n:=\mathbb{E}\left(\mathrm{e}^{t\sum_{i=1}^{n}X_i}\right)\cdot\mathrm{e}^{-nt(p+\varepsilon)}\le\exp\left[-nD_B(p+\varepsilon\parallel p)\right].
        $$

        Given the indepdency of $X_i$,
        $$
            \begin{aligned}
                g(t)^n &= \prod_{i=1}^n\mathbb{E}\left(\mathrm{e}^{tX_i}\right)\cdot\mathrm{e}^{-nt(p+\varepsilon)}\\
                &= (1-p+p\mathrm{e}^{t})^n\mathrm{e}^{-nt(p+\varepsilon)}\\
                &= \left[(1-p+p\mathrm{e}^t)\mathrm{e}^{-t(p+\varepsilon)}\right]^n\\
                &= \left[(1-p)\mathrm{e}^{-t(p+\varepsilon)}+p\mathrm{e}^{t(1-p-\varepsilon)}\right]^n.
            \end{aligned}
        $$

        We care about the minimum point of $g(t)$. Note that
        $$
            \begin{aligned}
                g'(t) &= -(1-p)(p+\varepsilon)\mathrm{e}^{-t(p+\varepsilon)}+p(1-p-\varepsilon)\mathrm{e}^{t(1-p-\varepsilon)}\\
                &= \left[-(1-p)(p+\varepsilon)+p(1-p-\varepsilon)\mathrm{e}^t\right]\mathrm{e}^{-t(p+\varepsilon)},
            \end{aligned}
        $$
        where the first item is a monotone increasing function of $t$ with the unique zero point (and  the minimum point)
        $$
            t_0=\ln\frac{(1-p)(p+\varepsilon)}{p(1-p-\varepsilon)}.
        $$

        It can be shown with patience that
        $$
            1-p+p\mathrm{e}^{t_0}=1-p+\frac{(1-p)(p+\varepsilon)}{1-p-\varepsilon}=\frac{1-p}{1-p-\varepsilon}
        $$
        and
        $$
            \mathrm{e}^{t_0(p+\varepsilon)}=\left[\frac{(1-p)(p+\varepsilon)}{p(1-p-\varepsilon)}\right]^{p+\varepsilon}.
        $$

        Finally,
        $$
            (1-p+p\mathrm{e}^{t_0})^{-1}\mathrm{e}^{t_0(p+\varepsilon)}=\left(\frac{p+\varepsilon}{p}\right)^{p+\varepsilon}\left(\frac{1-p-\varepsilon}{1-p}\right)^{1-p-\varepsilon}=\mathrm{e}^{D_B(p+\varepsilon\parallel p)}.
        $$
\end{itemize}
\end{answer}



\begin{problem}{5 (12')} ~
\begin{itemize}
    \item [(1)] (3') For what value of $d$ does the volume $V(d)$ of a $d$-dimensional unit ball take on its maximum? You don't need to prove your result.
    \item [(2)] (4') Consider drawing a random point $x$ from the unit ball in $\mathbb{R}^d$ (surface and interior) uniformly at random. What's the variance of $x_1$ (the first coordinate of $x$)? You don't need to prove your result.
    \item [(3)] (5') Recall the way we generate a random unit vector in $d$ dimensions. First, generate $d$ i.i.d samples $v_i\;(i\in [d])$ from a Gaussian distribution with $\mu=0$ and $\sigma^2=1$. Then, define the vector $v$ as $v=[v_1,v_2,...,v_d]$. Finally, return $\frac{v}{||v||_2}$, where $||v||_2=\sqrt{\sum_{i=1}^d v_i^2}$. Briefly explain why it generates a unit vector uniformly at random.
\end{itemize}
\end{problem}

\begin{answer} ~
\begin{itemize}
    \item [(1)] $d=5$.
    \item [(2)] $\mathrm{Var}(x_1)=\frac{1}{d+2}$.
    \item [(3)] For any unit vector $u$, $u\cdot v=\sum_{i=1}^d u_i v_i\sim \sum_{i=1}^d u_i\mathcal N(0,1)$, which is $\mathcal N(0,1)$ since $\sum_{i=1}^d u_i^2=1$. Thus, $u\cdot v$ has the same distribution for any unit vector $u$.
\end{itemize}
\end{answer}



\begin{problem}{6 (12')}
A 3-dimensional cube has vertices, edges and faces. In a $d$-dimensional cube, these components are called faces. A vertex is a $0$-dimensional face, an edge a $1$-dimensional face, etc. Answer the following problems. You don't need to prove your result.
\begin{itemize}
    \item [(1)] (3') For $0\leq k\leq d$, how many $k$-dimensional faces does a $d$-dimensional cube have?
    \item [(2)] (3') What is the total number of faces of all dimensions? The $d$-dimensional face is the cube itself which you can include in your count.
    \item [(3)] (3') What is the surface area of a unit cube in $d$-dimensions (a unit cube has side-length one in each dimension)?
    \item [(4)] (3') What is the surface area of the cube if the length of each side was 2?
\end{itemize}
\end{problem}

\begin{answer} ~
\begin{itemize}
    \item [(1)] $\binom{d}{k}2^{d-k}$.
    \item [(2)] $3^d$.
    \item [(3)] $2d$.
    \item [(4)] $2^dd$.
\end{itemize}
\end{answer}



\begin{problem}{7 (16')}
Consider the upper hemisphere of a unit-radius ball in $d$-dimensions. What is the height of the maximum volume cylinder that can be placed entirely inside the hemisphere?

\textit{[Hint: You need to consider all possible placement options.]}
\end{problem}

\begin{answer}
    For the cylinder whose height is not axis-aligned (which intersects with the bottom of the hemisphere at at most one single point), we can always rotate the hemisphere to keep the cylinder inside it and make the cylinder obviously not maximal [detailed discussion about 2 cases (whether the point is the origin) here is omitted since its geometrically clear enough]. Thus, we only need to consider axis-aligned cylinders:

    (a) For cylinders whose height, let to be $h\in(0,1)$, is vertical to the base of the hemisphere, its maximum volume is
    $$
        V_{\text{ver}}(h) = hV(d-1)(1-h^2)^{\frac{d-1}{2}}.
    $$
    Let
    $$
        V_{\text{ver}}'(h) = V(d-1)(1-dh^2)(1-h^2)^{\frac{d-3}{2}} = 0,
    $$
    we have argmax $h^*=\frac{1}{\sqrt{d}}$ and
    $$
        V_{\text{ver}}(h^*) = \frac{V(d-1)}{\sqrt{d}}\left(1-\frac{1}{d}\right)^{\frac{d-1}{2}}.
    $$

    (b) For cylinders whose height, let to be $2h\in(0,2)$, is horizontal to the base of the hemisphere, its maximum volume is
    $$
        V_{\text{hor}}(h) = 2hV(d-1)\left(\frac{1-h^2}{4}\right)^{\frac{d-1}{2}} \le V_{\text{ver}}(h).
    $$

    Therefore, $V_{\max}=V_{\text{ver}}(h^*)=\frac{V(d-1)}{\sqrt{d}}\left(1-\frac{1}{d}\right)^{\frac{d-1}{2}}$.
\end{answer}



\begin{problem}{8 (20')} 
Consider the following geometries in $d$-dimensional space:
\begin{align*}
    \Phi_d=\left\{(x_1,\cdots,x_d)\mid x_1^2+\cdots+x_d^2\leq 1\right\}\qquad
    \Omega_d=\left\{(x_1,\cdots,x_d)\mid |x_1|+\cdots+|x_d|\leq 1\right\}
\end{align*}
\begin{itemize}
    \item [(1)] (6') Consider the following two random processes:
    \begin{itemize}
        \item Pick two uniformly random unit vectors $x,y$ from the sphere in $d$ dimension.
        \item Pick a uniformly random plane passing through the origin in $d$ dimensions. Then pick two uniformly random unit vectors $w,z$ from the 2D sphere lying on that plane.
    \end{itemize}
    Determine whether $x$ and $w$ are identically distributed, and whether the pairs $(x,y)$ and $(w,z)$ are identically distributed. Briefly explain the reasons.
    \item [(2)] (6') Suppose $x$ and $y$ are sampled independently from $\Phi_d$. Denote random variable $Z=x^\top y$. Calculate $Var(Z)$.
    \item [(3)] (8') Calculate $V(\Omega_d)$ (the volume of $\Omega_d$) and $A(\Omega_d)$ (the surface area of $\Omega_d$).
\end{itemize}
\end{problem}

\begin{answer} ~
\begin{itemize}
    \item [(1)] The former case is true, because both $x$ and $w$ are uniformly sampled point on the surface of the unit ball. The later case is false for $d>2$, because $\mathbb{E}(|x\cdot y|)=\frac{1}{d}$ yet $\mathbb{E}(|w\cdot z|)=\frac{1}{2}$.
    \item [(2)] WLOG, let $y=(1,0,\cdots)$, and $Z=x\cdot y=x_1$. Given that
        $$
            \|x\|^2 = \sum_{i=1}^d x_i^2 \sim \rho(r):= d r^{d-1} (0\le r\le 1),
        $$
        and
        $$
            \mathbb{E}(x_i) = 0,\quad \mathrm{Var}(x_i)=\mathbb{E}(x_i^2),
        $$
        we have
        $$
            \begin{aligned}
                \mathrm{Var}(x_1) &= \mathbb{E}(x_1^2)\\
                &= \mathbb{E}\left(\frac{\|x\|^2}{d}\right)\\
                &= \int_0^1 r^2\frac{1}{d}d r^{d-1}\mathrm{d}r\\
                &= \frac{1}{d+2}.
            \end{aligned}
        $$
    \item [(3)] Integrating in one dimension, we have
        $$
            \begin{aligned}
                V(\Omega_d) &= \int_{-1}^1 V(\Omega_{d-1})(1-|x|)^{d-1}\mathrm{d}x\\
                &= 2V(\Omega_{d-1})\int_0^1(1-x)^{d-1}\mathrm{d}x\\
                &= \frac{2}{d}V(\Omega_{d-1}).
            \end{aligned}
        $$
        Given $V(\Omega_1)=2$, we get $V(\Omega_d)=\frac{2^d}{d!}$.

        As for $A(\Omega_d)$,
        $$
            \mathrm{d}\left(V(\Omega_d)r^d\right) = A(\Omega_d)r^{d-1}\frac{1}{\sqrt{d}}\mathrm{d}r
        $$
        so $A(\Omega_d)=\frac{2^d\sqrt{d}}{(d-1)!}$.
\end{itemize}
\end{answer}



\end{document}